% GENERATED FILE. Edit canonical lessons, then run npm run generate:books.
% canonical-source-hash: fnv1a64:629117cd978e70c1
% canonical-lessons: UR-C347-muamala, UR-C347-samjhauta, UR-C347-tanaza, UR-C347-ulajhna, UR-C347-chherna

\chapter{The Issue at Hand}
\label{ch:ur-the-issue-at-hand}
\hlchaptermodality{347}

\begin{chapteropening}
\textbf{By the end of this chapter:} \emph{I can say what a discussion is about and how it was settled.}

Complete the last lesson of chapter 347: I can say what a discussion is about and how it was settled.
\end{chapteropening}

\section[\texorpdfstring{\ur{معاملہ} (muāmalā)}{معاملہ (muāmalā)}]{\ur{معاملہ} (muāmalā) — a matter to be settled; an affair}

\label{lesson:UR-C347-muamala}



\begin{quote}
Before the new one: say the Urdu for welfare, then the Urdu for to make known.
\end{quote}

\subsection*{You'll want to know: \ur{معاملہ}}
\textbf{\ur{معاملہ}} — \emph{muāmalā} — \textquotedblleft{}a matter to be settled; an affair\textquotedblright{}.

\begin{quote}
\ur{یہ} \ur{معاملہ} \ur{بڑا} \ur{مشکل} \ur{ہے۔}
\end{quote}

\begin{quote}
\emph{yih muāmalā baṛā mushkil hai.}
\end{quote}

\begin{quote}
This is a very difficult affair.
\end{quote}

Masculine.

\begin{cousinweb}
\textbf{\ur{معاملہ}} (\emph{muāmalā}) is usually traced, through Persian, to Arabic \emph{muʿāmala}, \textquotedblleft{}dealing, a transaction\textquotedblright{}.
\end{cousinweb}

\begin{grammarlens}[title={What you've built}]
One more word for giving reasons and opinions.
\end{grammarlens}

\subsection*{Your turn}
\begin{itemize}
\raggedright
  \item \emph{Say it:} \emph{muāmalā}
  \item \emph{Say it:} \emph{muāmalā}, once more
  \item \emph{Say it:} \emph{yih muāmalā baṛā mushkil hai.} — This is a very difficult affair.
  \item \emph{Recall:} say the Urdu for welfare, then the Urdu for to make known, then say \emph{muāmalā} again
\end{itemize}

\begin{tcolorbox}[breakable,colback=teal!4,colframe=teal!35!black,title={Before you move on}]
What does \ur{معاملہ} mean? (\textquotedblleft{}A matter to be settled; an affair\textquotedblright{}.) Say it once more.
\end{tcolorbox}
\section[\texorpdfstring{\ur{سمجھوتا} (samjhautā)}{سمجھوتا (samjhautā)}]{\ur{سمجھوتا} (samjhautā) — a compromise; a settlement}

\label{lesson:UR-C347-samjhauta}



\begin{quote}
Before the new one: say the Urdu for to make known, then the Urdu for a matter to be settled.
\end{quote}

\subsection*{You'll want to know: \ur{سمجھوتا}}
\textbf{\ur{سمجھوتا}} — \emph{samjhautā} — \textquotedblleft{}a compromise; a settlement\textquotedblright{}.

\begin{quote}
\ur{دونوں} \ur{فریقوں} \ur{میں} \ur{سمجھوتا} \ur{ہو} \ur{گیا۔}
\end{quote}

\begin{quote}
\emph{donoṅ farīqoṅ meṅ samjhautā ho gayā.}
\end{quote}

\begin{quote}
The two sides reached a compromise.
\end{quote}

Masculine.

\begin{cousinweb}
\textbf{\ur{سمجھوتا}} (\emph{samjhautā}) is a noun made on \textbf{\ur{سمجھنا}}, \textquotedblleft{}to understand\textquotedblright{}, which is inherited: CDIAL lists the Hindi cousin \emph{samajhnā} under Sanskrit \emph{saṃbudhyate} (CDIAL 12959). An understanding both sides reach.
\end{cousinweb}

\begin{grammarlens}[title={What you've built}]
One more word for giving reasons and opinions.
\end{grammarlens}

\subsection*{Your turn}
\begin{itemize}
\raggedright
  \item \emph{Say it:} \emph{samjhautā}
  \item \emph{Say it:} \emph{samjhautā}, once more
  \item \emph{Say it:} \emph{donoṅ farīqoṅ meṅ samjhautā ho gayā.} — The two sides reached a compromise.
  \item \emph{Recall:} say the Urdu for to make known, then the Urdu for a matter to be settled, then say \emph{samjhautā} again
\end{itemize}

\begin{tcolorbox}[breakable,colback=teal!4,colframe=teal!35!black,title={Before you move on}]
What does \ur{سمجھوتا} mean? (\textquotedblleft{}A compromise; a settlement\textquotedblright{}.) Say it once more.
\end{tcolorbox}
\section[\texorpdfstring{\ur{تنازعہ} (tanāzā)}{تنازعہ (tanāzā)}]{\ur{تنازعہ} (tanāzā) — a dispute}

\label{lesson:UR-C347-tanaza}



\begin{quote}
Before the new one: say the Urdu for a matter to be settled, then the Urdu for a compromise.
\end{quote}

\subsection*{You'll want to know: \ur{تنازعہ}}
\textbf{\ur{تنازعہ}} — \emph{tanāzā} — \textquotedblleft{}a dispute\textquotedblright{}.

\begin{quote}
\ur{زمین} \ur{پر} \ur{دو} \ur{بھائیوں} \ur{میں} \ur{تنازعہ} \ur{تھا۔}
\end{quote}

\begin{quote}
\emph{zamīn par do bhāiyoṅ meṅ tanāzā thā.}
\end{quote}

\begin{quote}
Two brothers had a dispute over land.
\end{quote}

Masculine.

\begin{cousinweb}
\textbf{\ur{تنازعہ}} (\emph{tanāzā}) is usually traced, through Persian, to Arabic \emph{tanāzuʿ}, \textquotedblleft{}contending with one another\textquotedblright{}.
\end{cousinweb}

\begin{grammarlens}[title={What you've built}]
One more word for giving reasons and opinions.
\end{grammarlens}

\subsection*{Your turn}
\begin{itemize}
\raggedright
  \item \emph{Say it:} \emph{tanāzā}
  \item \emph{Say it:} \emph{tanāzā}, once more
  \item \emph{Say it:} \emph{zamīn par do bhāiyoṅ meṅ tanāzā thā.} — Two brothers had a dispute over land.
  \item \emph{Recall:} say the Urdu for a matter to be settled, then the Urdu for a compromise, then say \emph{tanāzā} again
\end{itemize}

\begin{tcolorbox}[breakable,colback=teal!4,colframe=teal!35!black,title={Before you move on}]
What does \ur{تنازعہ} mean? (\textquotedblleft{}A dispute\textquotedblright{}.) Say it once more.
\end{tcolorbox}
\section[\texorpdfstring{\ur{الجھنا} (ulajhnā)}{الجھنا (ulajhnā)}]{\ur{الجھنا} (ulajhnā) — to get tangled; to get bogged down}

\label{lesson:UR-C347-ulajhna}



\begin{quote}
Before the new one: say the Urdu for a compromise, then the Urdu for a dispute.
\end{quote}

\subsection*{You'll want to know: \ur{الجھنا}}
\textbf{\ur{الجھنا}} — \emph{ulajhnā} — \textquotedblleft{}to get tangled; to get bogged down\textquotedblright{}.

\begin{quote}
\ur{وہ} \ur{ایک} \ur{چھوٹی} \ur{بات} \ur{میں} \ur{الجھ} \ur{گئے۔}
\end{quote}

\begin{quote}
\emph{voh ek chhoṭī bāt meṅ ulajh gae.}
\end{quote}

\begin{quote}
They got bogged down in a small point.
\end{quote}

Past (it …): \ur{وہ} \ur{دھاگا} \ur{الجھ} \ur{گیا}, \emph{voh dhāgā ulajh gayā}.

\begin{cousinweb}
\textbf{\ur{الجھنا}} (\emph{ulajhnā}) is inherited: CDIAL lists the Hindi cousin \emph{ulajhnā} under Sanskrit \emph{uparundhati}, \textquotedblleft{}obstructs\textquotedblright{} (CDIAL 2221).
\end{cousinweb}

\begin{grammarlens}[title={What you've built}]
One more word for giving reasons and opinions.
\end{grammarlens}

\subsection*{Your turn}
\begin{itemize}
\raggedright
  \item \emph{Say it:} \emph{ulajhnā}
  \item \emph{Say it:} \emph{ulajhnā}, once more
  \item \emph{Say it:} \emph{voh ek chhoṭī bāt meṅ ulajh gae.} — They got bogged down in a small point.
  \item \emph{Recall:} say the Urdu for a compromise, then the Urdu for a dispute, then say \emph{ulajhnā} again
\end{itemize}

\begin{tcolorbox}[breakable,colback=teal!4,colframe=teal!35!black,title={Before you move on}]
What does \ur{الجھنا} mean? (\textquotedblleft{}To get tangled; to get bogged down\textquotedblright{}.) Say it once more.
\end{tcolorbox}
\section[\texorpdfstring{\ur{چھیڑنا} (chheṛnā)}{چھیڑنا (chheṛnā)}]{\ur{چھیڑنا} (chheṛnā) — to bring up a topic; to tease}

\label{lesson:UR-C347-chherna}



\begin{quote}
Before the new one: say the Urdu for a dispute, then the Urdu for to get tangled.
\end{quote}

\subsection*{You'll want to know: \ur{چھیڑنا}}
\textbf{\ur{چھیڑنا}} — \emph{chheṛnā} — \textquotedblleft{}to bring up a topic; to tease\textquotedblright{}.

\begin{quote}
\ur{اس} \ur{نے} \ur{میٹنگ} \ur{میں} \ur{پیسے} \ur{کی} \ur{بات} \ur{چھیڑی۔}
\end{quote}

\begin{quote}
\emph{us ne mīṭing meṅ paise kī bāt chheṛī.}
\end{quote}

\begin{quote}
She brought up the money question at the meeting.
\end{quote}

Past, with \ur{نے} (he or she …): \ur{اس} \ur{نے} \ur{نئی} \ur{بات} \ur{چھیڑی}, \emph{us ne naī bāt chheṛī}.

\begin{cousinweb}
\textbf{\ur{چھیڑنا}} (\emph{chheṛnā}) is inherited: CDIAL lists the Hindi cousin \emph{cheṛnā}, \textquotedblleft{}to touch, to irritate\textquotedblright{}, under \emph{*ciḍ}, \textquotedblleft{}be angry\textquotedblright{}, a form it reconstructs (CDIAL 4794).
\end{cousinweb}

\begin{grammarlens}[title={What you've built}]
One more word for giving reasons and opinions.
\end{grammarlens}

\subsection*{Your turn}
\begin{itemize}
\raggedright
  \item \emph{Say it:} \emph{chheṛnā}
  \item \emph{Say it:} \emph{chheṛnā}, once more
  \item \emph{Say it:} \emph{us ne mīṭing meṅ paise kī bāt chheṛī.} — She brought up the money question at the meeting.
  \item \emph{Recall:} say the Urdu for a dispute, then the Urdu for to get tangled, then say \emph{chheṛnā} again
\end{itemize}

\begin{tcolorbox}[breakable,colback=teal!4,colframe=teal!35!black,title={Before you move on}]
What does \ur{چھیڑنا} mean? (\textquotedblleft{}To bring up a topic; to tease\textquotedblright{}.) Say it once more.
\end{tcolorbox}
